\section{实现}
\label{app:implementation}

\subsection{条件 token}
\label{app:tokens}
\paragraph{当前图像。}
冻结的 \cfgVisualEncoder{} \cite{mae} 将 $224\times224$ agent 视角 RGB 编码为 \cfgPatches{} 个图像块 token。我们在开发场景上比较冻结的 DINOv3 编码器 \cite{dinov3} 和与策略共同训练的 CNN 后选择它（附录~\ref{app:design_shared}）。每个图像块接收视角嵌入、有效度量深度均值的线性嵌入，以及反投影到基座坐标系的位置的线性嵌入。该位置根据固定源像素与相机内外参计算，也用于图像查询、键及共享相同源点的场景查询的三维旋转位置嵌入。目标、机器人和动作 token 不使用旋转位置嵌入。仅训练时，条件 token 上的线性解码器重建当前目标框内的 RGB（在 $\mathcal L_{\rm aux}$ 中权重为 0.1），从而在不增加输入的情况下把图像块特征与目标关联。

\paragraph{目标运动历史。}
因果跟踪器（附录~\ref{app:tracker}）在采集时刻输出位姿。每次查询时，在 $-350,-300,\ldots,0$\,ms 偏移处对已输出轨迹重采样，在估计之间对平移线性插值、对旋转沿测地线插值。不进行外推，也不跨越超过两个相机周期的缺口插值；两侧没有估计的偏移标记为无效。每个运动 token 为 25 维描述符
$(\vect p,\log R,\Delta\vect p,\log\Delta R,\Delta t,\vect v,\boldsymbol\omega,\vect a,\boldsymbol\alpha)$，
各部分维数为 $(3,3,3,3,1,3,3,3,3)$，共 \cfgHistoryLen{} 个 token。槽位 0--6 描述相邻偏移间的七个区间，槽位 7 保存当前绝对位姿。速度由区间位移除以经过时间得到，加速度由相邻速度作差得到。各维按训练统计标准化，无效维置零，再附加三项有效性标记（位姿、区间、加速度），送入带时间嵌入的 $28\to\cfgWidth\to\cfgWidth$ GELU MLP。跟踪器报告目标丢失时，用可学习嵌入替代位姿。

\paragraph{机器人状态。}
同一组偏移处的 \cfgHistoryLen{} 个末端执行器状态（位置、旋转向量、实测夹爪开度）构成机器人 token。另一个 token 保存首个生效时刻 $\tapp_{t,0}$ 的末端指令位姿，作为动作块增量的参考（附录~\ref{app:action}）。

\subsection{位姿跟踪器}
\label{app:tracker}
SAM~2.1（Hiera-small）从首帧框因果传播目标掩码 \cite{sam2}。仿真中，框来自目标的渲染掩码；真机上，框由 Grounding DINO 根据指令中的物体短语给出 \cite{groundingdino}。FoundationPose 使用发布权重，在全分辨率 RGB-D 帧（\fillin{$640\times480$}）中配准并跟踪物体 \cite{foundationpose}。仿真中，物体模型为目标的仿真器网格；真机上为 CAD 模型或每个物体一次性采集的 \fillin{16} 个带掩码参考视图。跟踪器在独立线程中运行，每个估计附有采集时刻、可用时刻和有效性标记。配准失败、掩码范围内深度丢失，或帧间跳变超过 \fillin{X}\,mm 时，标记为无效。

单目三维检测器位姿（第~\mainref{sec:main_results}~节）无需物体模型。对每个带掩码的帧，WildDet3D \cite{wilddet3d} 从 RGB 帧、紧致二维框及内参预测所提示目标的度量三维框。框轴存在由对称性导致的歧义，因此如 DROID 处理（附录~\ref{app:droid_target}），我们从 24 种轴排列与符号翻转中选择最接近前帧朝向的一种，把每个框映射到该 episode 内统一的物体坐标系。旋转对称目标的估计、标签与误差均投影去除绕对称轴的旋转。表~\ref{tab:supp_tracker} 报告两种跟踪器的精度与输入分辨率的影响。

% Planned results: selected marker-free pipeline; real errors require an independent marker-free reference with stated uncertainty. No tracker ranking or ball rotation experiment.
\begin{table}[t]
\centering\footnotesize
\caption{\tablelang{\textbf{Real-robot pose tracking.} Median/p95 translation (mm) and rotation ($^\circ$) error of the tracker against a marker-free reference, the uncertainty of that reference, and the fraction of lost observations. Ball rotation is not used.}{\textbf{真机位姿跟踪。}跟踪器相对无标记参照的平移（mm）与旋转（$^\circ$）误差的中位数/p95、该参照的不确定度，以及观测丢失比例。球体不使用旋转。}}
\label{tab:supp_tracker}
\setlength{\tabcolsep}{2.8pt}
\begin{tabular}{@{}lcccc@{}}
\toprule
\tablelang{Target}{目标} & mm & $^\circ$ & \tablelang{Ref.\ uncertainty}{参照不确定度} & \tablelang{Lost}{丢失} \\
\midrule
\tablelang{Real rigid object}{真机刚体} & \resultplaceholder[{0.0/0.0}] & \resultplaceholder[{0.0/0.0}] & \resultplaceholder & \resultplaceholder \\
\tablelang{Real ball}{真机球体} & \resultplaceholder[{0.0/0.0}] & -- & \resultplaceholder & \resultplaceholder \\
\bottomrule
\end{tabular}
\end{table}


\subsection{世界流}
\label{app:world}
世界由一个可学习物体查询（$\mathcal Q$）和每个网格点一个场景查询（$\mathcal F$）组成，宽度均为 \cfgWidth{}。它们由可学习嵌入和位置编码初始化，关注条件及彼此，并采用确定性回归，因此不输入未来标签或噪声。读出头嵌入每个时域的物理偏移 $\delta_h$，从而可在其他偏移监督同一组查询，预训练即采用此方式（附录~\ref{app:droid}）。

\paragraph{目标。}
MLP 从物体查询的最终状态读出轨迹 $\{(\vect d_h,\boldsymbol\phi_h)\}_{h=1}^{\cfgHorizons}$。对有效时域 $\mathcal V$，
\begin{equation}
\mathcal L_P=\frac1{|\mathcal V|}\sum_{h\in\mathcal V}\Big[\|\hat{\vect d}_h-\vect d_h\|^2
+\lambda_R\,d_{\mathrm{SO}(3)}\big(\hat R_h,R_h\big)^2\Big],
\end{equation}
其中 $\lambda_R=0.1$，$\hat R_h=\Exp(\hat{\boldsymbol\phi}_h^\wedge)R(\tau_t)$。对称目标的 $d_{\mathrm{SO}(3)}$ 忽略绕对称轴旋转。我们使用左相对旋转向量与累积平移，而非会通过雅可比耦合平移与旋转的 twist 指数映射。预测为条件均值，因此会平均多模态未来，例如碰撞发生前的多种可能结果。

\paragraph{场景流。}
流表示 $N=\cfgContextPoints$ 个源点在 \cfgHorizons{} 个时域上的累积三维位移 $F\in\R^{N\times\cfgHorizons\times3}$，位于基座坐标系，并附有可见性 $o\in\{0,1\}^{N\times\cfgHorizons}$。每个场景查询包含源位置与角色嵌入。角色（目标、机器人或环境）来自跟踪器掩码以及机器人的渲染或运动学掩码。线性读出给出每点各时域的位移与可见性 logit。位移除以在训练标签上拟合的固定尺度，损失为
\begin{equation}
\mathcal L_F=\frac1{|\mathcal V_F|}\sum_{(i,h)\in\mathcal V_F}o_{ih}\big\|\hat F_{ih}-F_{ih}\big\|^2
+\mathrm{BCE}(\hat o,o),
\end{equation}
其中 $\mathcal V_F$ 包含有效标签，因此被遮挡点仅监督可见性。

\paragraph{外观目标。}
场景查询最终状态上的图像块解码器，在各时域预测三类目标：该块的度量深度（有效深度上的平方误差）、冻结 DINOv3 ViT-L/16 教师 \cite{dinov3} 在 $14\times14$ 网格上的特征（逐通道层归一化后的平方误差），以及运动表面的 RGB（位移超过 1\,mm 的表面，包含机器人，在该掩码内计算平方误差）。各目标在 $\mathcal L_{\rm aux}$ 中权重均为 0.1。部署时跳过这些头，没有 token 读取其输出，因此它们仅塑造场景查询状态。

\subsection{动作流}
\label{app:action}
每个动作块包含 \cfgChunk{} 条七维指令：基座坐标系平移增量、左相对旋转向量和夹爪指令。指令 $j$ 相对 $\tapp_{t,j-1}$ 的指令位姿定义增量；当 $j=0$ 时，参考位姿为仍在执行的上一动作块给 $\tapp_{t,0}$ 指定的位姿。因此动作块锚定于控制器精确已知的位姿，而非采集时刻机械臂的实测位姿。每个动作 token 加入其领先时长 $\Delta_{t,j}=\ell_t+j\Delta_a$ 的正弦嵌入。训练时每个窗口采样延迟 $\ell\sim\mathcal U[0,\cfgLatencyMax]$，指令 $j$ 的标签为示范中 $\tau_t+\ell+j\Delta_a$ 时刻的指令增量，在控制步之间插值。归一化使用训练划分中有效、未重复的已发出指令的均值与总体标准差。流匹配在该坐标下使用单位高斯噪声，每块一个 $u$，动作 token 通过自适应层归一化接收 $u$。部署时逆归一化、解码旋转向量，并逐轴将平移和旋转截断到 0.10\,m 和 0.5\,rad。夹爪按指令符号动作。

\subsection{耦合}
\label{app:coupling}
注意力掩码作用于条件（$\mathcal Z$）、物体查询（$\mathcal Q$）、场景查询（$\mathcal F$）和动作（$\mathcal A$）四组 token。条件 token 仅关注 $\mathcal Z$；世界查询关注 $\mathcal Z$、$\mathcal Q$、$\mathcal F$；动作 token 关注 $\mathcal Z$、$\mathcal A$ 以及耦合增加的连接（附录~\ref{app:design_act}）。阻断边在训练和部署时的每层、每头 logit 均为 $-\infty$。因此，动作不关注的世界查询不能影响动作，未来标签也不能经条件流入动作。所有耦合共享模块、初始化与损失，仅掩码不同，R5 则另有几何编码器。两流之间不截断梯度。除非打开动作到世界的边，世界既不关注动作，也不关注流时间，因此 R1、R2、R5 的世界状态每次查询只计算一次，在 Euler 步间复用；R3、R4 则每步重算。
\Qthree{1}{选定的 R2 耦合中，动作 token 在每层关注由该层输入状态计算的 $\mathcal Q$、$\mathcal F$ 的键和值。部署时计算一次世界，再以 \cfgEuler{} 步 Euler 积分生成动作。}%
\Qthree{2}{选定的 R5 耦合中，解码后的目标轨迹、场景查询的位移与可见性概率，以及连续 6D 形式的旋转，经 0.37M 参数的几何编码器形成每查询一个几何 token。动作除关注 $\mathcal Z$ 外也关注这些 token，不读取任何世界隐藏状态。训练使用预测几何，动作损失经其回传，不采用教师强制。部署时预测一次世界，跳过外观头，再以 \cfgEuler{} 步 Euler 积分生成动作。}%
\Qthree{3}{选定的 R1 耦合没有打开从世界到动作的边。世界损失塑造共享参数，部署时只运行动作流，执行 \cfgEuler{} 步 Euler 积分。}
\Bbone{1}{主干各组共享注意力投影，世界查询与动作 token 使用独立前馈网络（扩展比 4）。世界查询采用层归一化，动作 token 采用自适应层归一化，可训练主干参数为 \cfgBackboneParams{}。}%
\Bbone{2}{主干具有独立的世界与动作参数组，各自包含注意力投影、自适应层归一化和前馈网络（扩展比 3）。条件 token 由世界组处理，因此世界损失也塑造动作读取的条件表示。}
附录~\ref{app:design_act} 的另一种参数化，其可训练参数量相差 0.013\%。

\subsection{训练配方}
\label{app:recipe}
表~\ref{tab:hyper} 列出 \bench 训练配方。每条示范提供查询间隔为 \fillin{8} 个控制步的窗口，约每条示范 \fillin{14} 个，每窗独立采样延迟。先导试验中，全局 batch 64 配合两倍学习率使训练快约 \fillin{2.5} 倍，但跨种子稳定性较低，因此保留 16。一次训练在 \cfgTrainHardwareZh{} 上约需 \cfgTrainHours{} 小时。DROID 基座与受控对照采用附录~\ref{app:scale} 的配方。

\begin{table}[h]
\centering\small
\caption{\bench 上的训练超参数。}
\label{tab:hyper}
\begin{tabularx}{\columnwidth}{@{}lX@{}}
\toprule
项目 & 取值 \\
\midrule
主干 & \cfgLayers{} 层，宽度 \cfgWidth{}，\cfgHeads{} 个头，\cfgBackboneParams{} 参数 \\
编码器 & 冻结 \cfgVisualEncoder{}，$224\times224$ 输入 \\
优化器 & AdamW，$\beta=(0.9,0.95)$，$\epsilon=10^{-8}$，权重衰减 0.05 \\
学习率 & $2\times10^{-4}$，1{,}000 次预热更新，余弦衰减到 0.1 \\
batch 与更新 & 全局 batch 16，4 个 epoch（40{,}104 次更新，641{,}664 个窗口） \\
损失权重 & $\lambda_P=\lambda_F=1$，$\lambda_R=0.1$，各目标 $\lambda_{\rm aux}=0.1$ \\
延迟采样 & 每窗 $\ell\sim\mathcal U[0,\cfgLatencyMax]$ \\
精度 & BF16 配合 TF32，梯度裁剪阈值 1，无 dropout \\
采样 & 动作使用 \cfgEuler{} 步 Euler 积分 \\
推理 & FP16、编译、条件 token 缓存 \\
硬件 & \cfgTrainHardwareZh{}，每次约 \cfgTrainHours{} 小时 \\
\bottomrule
\end{tabularx}
\end{table}

\subsection{部署}
\label{app:deployment}
每次查询编码最新观测，在查询内复用缓存的条件 token，并发出按各自生效时刻调度的动作块（式~\mainref{eq:apptime}）。查询前，执行器将 $\ell_t$ 设为延迟预算，即最近 50 次实测推理时间的 p95 加平台执行延迟。若推理提前结束，指令等待至预定时刻；若推理超时，丢弃已过时指令，执行其余指令。附录~\ref{app:sweeps} 说明所有方法共享的执行器，附录~\ref{app:timing} 说明计时。
